\begin{lemma}[Appending a terminal maximum preserves the maximum property]
Let \(T\) be an extended Tashkinov tree with positive rung number \(n\),
certified series \(\mathcal T\), maximum property in the sense of
\noderef{MaximumProperty}, and a terminal maximum closure in the sense of
\noderef{TerminalClosureMaximum}.  Append one further certified
Algorithm~3.1 step to \(\mathcal T\): more precisely, let
\(\mathcal T^+\) be a certified Tashkinov series with \(n+1\) noninitial
iterations whose complete indexed data is exactly the old data followed by
one supplied last tree, coloring, stored set, exact connecting-edge set,
extension kind, chosen and outside vertices, and two colors, as defined by
\noderef{TashkinovSeriesDataAppend}.

Then there exists an extended Tashkinov tree displayed on exactly the final
tree of \(\mathcal T^+\), with generating coloring exactly the terminal
coloring of \(\mathcal T^+\), and this ETT has the maximum property.  In
particular, its displayed tree and generating coloring are the entries at
indices \(n+2\) and \(n+1\), respectively.
\end{lemma}
\begin{proof}
Because \(n+1>0\), apply
\noderef{TashkinovSeriesFinalETTExists} to the certified appended series
\(\mathcal T^+\).  It gives an ETT \(T^+\) displayed on exactly the last
tree of \(\mathcal T^+\), with exactly its terminal coloring, rung number
\(n+1\), and dependent certified series exactly \(\mathcal T^+\).  It
therefore remains to verify the two clauses of
\noderef{MaximumProperty} for this ETT.

For the global first-tree clause, the index \(1\) is an old tree index.
The exact append equations in \noderef{TashkinovSeriesDataAppend} give
\(T_1^+=T_1\) as ordered tree-sequences.  The maximum property of the old
ETT says \(|V(T_1)|\) equals the Tashkinov order of \(G\).  Rewriting along
this equality proves the first clause for \(T^+\).

For the stagewise clause, fix an integer \(i\) with
\(1\leq i<n+1\).  Let \(\pi\) be any proper partial coloring leaving
exactly \(e_0\) uncolored and stable for the stage-\(i\) tree, computed
stored set, and reference coloring in \(\mathcal T^+\).  Let \(U\) be a
seed satisfying its PE or SE/RE route, and let \(U'\) be a closure of \(U\)
under \(\pi\).  We must prove
\[
                         |V(U')|\leq |V(T_{i+1}^+)|.       \tag{1}
\]

First suppose \(i<n\).  Both \(i\) and \(i+1\) are old indices.  By the
array equations in \noderef{TashkinovSeriesDataAppend}, the tree \(T_i^+\),
reference coloring \(\varphi_i^+\), kind \(\Theta_i^+\), exact
connecting-edge set \(F_i^+\), and comparison tree \(T_{i+1}^+\) equal
their old counterparts.  The stored set also agrees exactly:
\noderef{TashkinovStoredColorsAppend} gives \(D_i^+=D_i\).  Consequently
the assumed stability of \(\pi\) is exactly
\((T_i,D_i,\varphi_i)\)-stability.

The prescribed seed route transports without loss of edge identity.  In
the PE case, \(\Theta_i^+=PE\) and \(U=T_i^+\) rewrite to
\(\Theta_i=PE\) and \(U=T_i\).  In the SE/RE case, the same equalities
rewrite the alternative for the kind, the exact singleton equation
\(F_i^+=\{f\}\), and the assertion that \(U=T_i^++f\) to the corresponding
old assertions \(F_i=\{f\}\) and \(U=T_i+f\).  Thus every hypothesis of
the old stage-\(i\) clause of \noderef{MaximumProperty} holds for the same
\(\pi,U,U'\).  That clause yields
\(|V(U')|\leq|V(T_{i+1})|\), and the append equality
\(T_{i+1}^+=T_{i+1}\) turns this into (1).

It remains to consider \(i\not<n\).  Together with \(i<n+1\), ordinary
integer order gives \(i=n\).  The tree \(T_n^+\), reference coloring
\(\varphi_n^+\), kind \(\Theta_n^+\), and exact connecting-edge set
\(F_n^+\) are again old entries by
\noderef{TashkinovSeriesDataAppend}, and
\noderef{TashkinovStoredColorsAppend} gives \(D_n^+=D_n\).  Exactly the
same rewriting as in the preceding paragraph turns the stability and seed
route assumptions into the terminal assumptions for the old series.
The terminal maximum-closure hypothesis
\noderef{TerminalClosureMaximum} therefore gives
\[
                         |V(U')|\leq |V(T_{n+1})|.         \tag{2}
\]
The comparison tree required here is \(T_{i+1}^+=T_{n+1}^+\), not the
newly supplied last tree \(T_{n+2}^+\).  Since the old final tree at index
\(n+1\) is retained by the append operation, the exact append equation is
\(T_{n+1}^+=T_{n+1}\).  Rewriting (2) proves (1) in the remaining case.

Every index \(1\leq i<n+1\) has now been covered.  Hence the global
first-tree equation and all stagewise bounds in
\noderef{MaximumProperty} hold for the exact-final ETT \(T^+\), as
claimed.
\end{proof}
